\BeleskeID{02-s036}
% element: 02-s036:e01
% element: 02-s036:e02
% element: 02-s036:d01

Indeks u karti računamo iz kombinacije \(DCBA\) kao \(8D+4C+2B+A\). Grejov redosled na obe ose obezbeđuje promenu jednog bita pri prelazu preko zajedničke stranice. Preko leve i desne ivice menja se samo \(B\), a preko gornje i donje samo \(D\).

Ako kartu zamislimo presavijenu oko vertikalne ose, polja preko leve i desne ivice dobijaju susedne stranice. Analogno, presavijanje oko horizontalne ose prikazuje susedstvo preko gornje i donje ivice. Ovaj geometrijski prikaz pomaže da uočimo iste parove koje određuje razlika u jednom bitu.

Četiri ugaona polja 0, 2, 8 i 10 zato mogu zajedno činiti pravilnu grupu od četiri polja, iako su na ravnom crtežu razdvojena. Za njih važi \(C=0\) i \(A=0\), dok se \(D\) i \(B\) menjaju. Ako su sva četiri polja jedinice, grupi odgovara proizvod \(\bar C\bar A\). Dijagonalna polja bez takvog susedstva ne grupišemo samo zato što se dodiruju uglom.
